\SourcePage{611}{614}
\section{Analysis of the general formula}

The formulas obtained above apply, in principle, to scattering by any field
$U(r)$ that tends to zero at infinity.  Their analysis therefore reduces to
examining the properties of the phases $\delta_l$ which occur in them.

To estimate the order of magnitude of $\delta_l$ at large $l$, we use the
fact that the motion is quasiclassical for large $l$ (see \S~49).  The phase
of the wave function is consequently determined by the integral
\[
 \int_{r_0}^{r}\sqrt{k^2-\frac{(l+1/2)^2}{r^2}
 -\frac{2mU(r)}{\hbar^2}}\,dr+\frac\pi4,
\]
where $r_0$ is a root of the expression under the radical (the region
$r>r_0$ is classically accessible).  From this we subtract the phase
\[
 \int_{r_0}^{r}\sqrt{k^2-\frac{(l+1/2)^2}{r^2}}\,dr+\frac\pi4
\]
of the free-motion wave function and then let $r\to\infty$; by definition,
the result is $\delta_l$.  For large $l$, the value of $r_0$ is also large.
Hence $U(r)$ is small throughout the integration region, and approximately

\SourcePage{612}{615}
\begin{equation}
 \delta_l=-\int_{r_0}^{\infty}
 \frac{mU(r)\,dr}{\hbar^2\sqrt{k^2-(l+1/2)^2/r^2}}.
 \tag{124.1}
\end{equation}
In order of magnitude, this integral (provided it converges) is
\begin{equation}
 \delta_l\sim\frac{mU(r_0)r_0}{k\hbar^2}.
 \tag{124.2}
\end{equation}
The corresponding magnitude of $r_0$ is $r_0\sim l/k$.

If $U(r)$ approaches zero at infinity as $r^{-n}$ with $n>1$, the integral
(124.1) converges and the phases $\delta_l$ are finite.  For $n\leqslant1$,
on the contrary, the integral diverges, and the phases $\delta_l$ are
infinite.  This statement applies to every $l$: convergence of (124.1)
depends on the large-$r$ behavior of $U(r)$, while at large distances (where
the field $U(r)$ is already weak) radial motion is quasiclassical for any
$l$.  The interpretation of formulas (123.11) and (123.12) when $\delta_l$
is infinite will be given below.

First consider convergence of the series (123.12), which represents the
total scattering cross-section.  At large $l$ the phases satisfy
$\delta_l\ll1$, as follows from (124.1) when $U(r)$ decreases faster than
$1/r$.  We may therefore put $\sin^2\delta_l\approx\delta_l^2$, so that the
tail of the series (123.12) has the order of
$\sum_{l\gg1}l\delta_l^2$.  The usual integral test shows that this series
converges if $\int_0^\infty l\delta_l^2\,dl$ converges.  Substitution of
(124.2), followed by replacing $l$ with $kr_0$, gives the integral
\[
 \int^{\infty}U^2(r_0)r_0^3\,dr_0.
\]
When $U(r)$ falls off at infinity as $r^{-n}$ with $n>2$, this integral
converges and the total cross-section is finite.  If, however, the field
$U(r)$ decreases as $1/r^2$ or more slowly, the total cross-section is
infinite.  Physically, this is because a slowly decreasing field gives a
very large probability of scattering through small angles.  Recall in this
connection that, in classical mechanics, a particle passing at any finite
impact parameter $\rho$, however large, still undergoes a deflection through
some small

\SourcePage{613}{616}
but nonzero angle whenever the field vanishes only as $r\to\infty$; the total
scattering cross-section is therefore infinite for every law by which $U(r)$
decreases.\footnote{This appears as the divergence of the integral
$\int2\pi\rho\,d\rho$ which defines the total cross-section in classical
mechanics.}  This argument does not apply in quantum mechanics, if only
because scattering through a given angle can be discussed only when that
angle is large compared with the uncertainty in the direction of motion.  If
the impact parameter is known to an accuracy $\Delta\rho$, it produces an
uncertainty $\sim\hbar/\Delta\rho$ in the transverse momentum component, and
thus an angular uncertainty $\sim\hbar/(mv\Delta\rho)$.

Because small-angle scattering is so important when $U(r)$ decreases slowly,
one naturally asks whether the scattering amplitude $f(\theta)$ might diverge
at $\theta=0$ even when $U(r)$ falls off faster than $1/r^2$.  On setting
$\theta=0$ in (123.11), the distant terms of the sum give an expression
proportional to $\sum_{l\gg1}l\delta_l$.  The same reasoning as before reduces
the condition for finiteness of the sum to convergence of
\[
 \int^{\infty}U(r_0)r_0^2\,dr_0,
\]
which already diverges for $U(r)\sim r^{-n}$ when $n\leqslant3$.  Thus the
scattering amplitude becomes infinite at $\theta=0$ ($n\leqslant1$) in fields
which fall off as $1/r^3$ or more slowly.

Finally, consider the case in which the phase $\delta_l$ itself is infinite,
as happens for $U(r)\sim r^{-n}$ with $n\leqslant1$.  The preceding results
make it clear in advance that, with so slow a decrease of the field, both the
total cross-section and the scattering amplitude at $\theta=0$ are infinite.
It remains to determine $f(\theta)$ for $\theta\ne0$.  We first note the
formula\footnote{This formula is the expansion of the $\delta$-function in
Legendre polynomials and follows directly on multiplying both sides by
$\sin\theta P_l(\cos\theta)$ and integrating with respect to $d\theta$.  Here
the integral $\int_0^\infty\delta(x)\,dx$ of the even function $\delta(x)$ is
taken to be $1/2$.}
\begin{equation}
 \sum_{l=0}^{\infty}(2l+1)P_l(\cos\theta)=4\delta(1-\cos\theta).
 \tag{124.3}
\end{equation}
In other words, this sum vanishes for every $\theta\ne0$.  Consequently, for
$\theta\ne0$ the unity in square brackets in every term of (123.11) may be

\SourcePage{614}{617}
omitted, leaving
\begin{equation}
 f(\theta)=\frac1{2ik}\sum_{l=0}^{\infty}(2l+1)
 P_l(\cos\theta)e^{2i\delta_l}.
 \tag{124.4}
\end{equation}
Multiplication of the right-hand side by the constant factor
$\exp(-2i\delta_0)$ does not affect the cross-section, which is determined by
$|f(\theta)|^2$; it merely changes the phase of the complex function
$f(\theta)$ by an inessential constant.  On the other hand, in the difference
$\delta_l-\delta_0$ between expressions (124.1), the divergent integral of
$U(r)$ cancels and a finite quantity remains.  Thus, in the case considered,
the scattering amplitude may be calculated from
\begin{equation}
 f(\theta)=\frac1{2ik}\sum_{l=0}^{\infty}(2l+1)
 P_l(\cos\theta)e^{2i(\delta_l-\delta_0)}.
 \tag{124.5}
\end{equation}
